% ============================================================
% 分片：s03g  对应原文 PDF 第 55--60 页（含第 55 页顶部 §3.2.6 收尾三段补遗）
% 负责章节：3.3 Holomorphic and mixed topological-holomorphic descent；
%           3.3.1 Holomorphic twists in 2d；3.3.2 Holomorphic-topological twists in 3d；
%           3.3.3 Homotopy actions of vertex algebras and Poisson vertex algebras
% 译者：translator-01   审校状态：初稿
% ============================================================

% [边界补遗: 原文第 55 页顶部三段属 §3.2.6 收尾；因分片按页切分，故并入本分片，勿与 s03f 重复]
% [跨页接续: 原文第 54 页/55 页]
% ==== 原文 PDF 第 55 页 ====
更一般地，若用 BV 形式表述任何规范理论，所得作用量的一阶形式自然导向一个（$L_\infty$ 型的）Maurer--Cartan（MC）方程 [100]。
由（经典）BV 主方程得到的 MC 方程可以给出形变理论的解释：它涉及一个带有 BV 微分的底层形式流形；$L_\infty$ 代数相当普遍地已渗透到微扰量子场论与弦论之中！
这一视角还与遵循 Costello 与 Gwilliam [55] 的、把量子场论表述为分解代数的现代框架完美契合，细节请参阅该文献。

我们还看到，扭变 QFT 中的后代可以作为形变加入作用量。
在那里，（同伦）代数的 MC 方程同样在幕后起作用。
在这一语境中，形变元素 $x$ 应视为其后代使作用量发生形变的那个算子，而 $d_A$ 由扭变后的 BRST 微分给出。
$A$ 是扭变 QFT 的算子代数。
例如，在我们此前大部分时间都在讨论的 $A$ 型与 $B$ 型扭变中，它们分别是相应的拓扑手征环与扭手征环。
$x$ 是 MC 元素这一事实，等价于形变项在任意阶微扰论中都是 $d_A$ 不变的（物理上，即规范不变的）。
这可以通过直接的路径积分运算证明；DG 结合代数语境下的计算见 [87, 101]，$A_\infty$ 情形下的表达式见 [102]。
我们要强调：此处 MC 方程不是经典运动方程，而是在任意阶微扰论中描述量子层面的规范不变性，它往往导致高阶代数结构的丰富形变。
% [Tutor Note: 物理图像：这里 $x$ 的 MC 方程与 (3.59) 形式相同，但含义不同——它并非运动方程，而是「任意阶微扰下形变项仍为 $d_A$ 闭」的代数编码；
% 同伦代数把这一条件提升为 $L_\infty$ 或 $A_\infty$ 的 MC 方程。]

人们可能还想研究由把两个原本不耦合的部分耦合起来而产生的微扰形变，例如带某种局域序型缺陷（order-type defect）的（扭变）体 QFT，或开--闭（拓扑）弦理论。
这导致同调代数中一个称为 Koszul 对偶的概念出现，它刻画了所得形变在任意阶微扰论中的 (BV)-BRST 不变性。
综述与原始文献见 [87]。
由于在研究两个算子代数的微扰形变（耦合）时 Koszul 对偶普遍出现，它在开--闭弦理论中的出现对全息有重要应用——全息可视为一种开--闭弦对偶 [101, 103]。
这些联系仍在活跃的研究与发展之中。

% [说明: 以上三段为 §3.2.6 收尾（原文第 55 页上半）；以下自 §3.3 节标题起。]
\subsection{全纯下降与混合拓扑--全纯下降}

我们此前一直关注拓扑扭变理论中的拓扑下降。在早先关于扭变的讨论中我们注意到，一般而言，人们可能反而会得到沿某些或全部扭变方向的全纯依赖性。
下降程序有一个全纯版本，而扭变理论中由此得到的代数结构的丰富化是一个活跃的研究领域。
如上所述，拓扑理论中支撑下降与碰撞的代数结构是一个 $E_d$ 代数。
对于全纯或混合全纯–拓扑情形，相应的结构尚未被很好地理解。

% [Tutor Note: 物理图像：拓扑扭变（$n=d$）给出 $E_d$ 代数；完全全纯扭变（$n=0$）给出的是（高阶）顶点/手性代数；混合情形（$0<n<d$）介于两者之间，
% 其代数结构正是本节要构造的对象。原文指出后者的完整代数刻画仍是开放问题。]
扭变在不同维度上的结构意味着：偶数维中的极小扭变将是完全全纯的（$n=0$），而奇数维中的极小扭变将有一个拓扑方向（$n=1$）[15, 16]。
稍更一般地，扭变超荷 $Q$ 将平凡化比如说 $n$ 个实平移与 $m$ 个反全纯平移，而时空维数为 $d=n+2m$。
若取局部实坐标 $x^\mu$（$\mu=1,\dots,n$）与复坐标 $z^a,\bar z^{\bar a}$（$a,\bar a=1,\dots,m$），我们预期找到费米型算子 $Q_\mu,Q_{\bar a}$，它们与动量对易，%
% ==== 原文 PDF 第 56 页 ====
% [跨页接续: 原文第55页/56页，该句自上一页续接]
并且彼此反对易，使得
\begin{equation}\tag{3.60}
\{Q,Q_\mu\}=iP_\mu \qquad \{Q,Q_{\bar a}\}=iP_{\bar a}.
\end{equation}

与纯粹拓扑下降中一样，可以用这些算子构造局域算子 $O$ 的混合全纯–拓扑后代 $O^{(k)}$，它们满足混合全纯–拓扑下降方程
\begin{equation}\tag{3.61}
QO^{(k)}=d'O^{(k-1)}+(QO)^{(k)},
\end{equation}
其中局部地 $d'=dx^\mu\partial_\mu+dz^{\bar a}\partial_{\bar a}=d_{\mathbb{R}^n}+\bar\partial_{\mathbb{C}^m}$ 是 $\mathbb{R}^n$ 上的 de Rham 微分与 $\mathbb{C}^m$ 上的 Dolbeault 微分之和。
% [原文疑误: 此处按 Dolbeault 微分应为 $d\bar z^{\bar a}\partial_{\bar a}$，原文 $dz^{\bar a}$ 的 $z$ 上漏印横杠（依据：紧接的等式 $d_{\mathbb{R}^n}+\bar\partial_{\mathbb{C}^m}$ 以及 $P_{\bar a}$ 的定义）；译文照抄原样。]
正如在完全拓扑的设定中一样，$O^{(k)}$ 显式地包含 $Q$ 闭局域算子对 $x^\mu$ 与 $\bar z^{\bar a}$ 的依赖信息。
特别地，在 $Q$ 上同调层面可推出：$Q$ 闭局域算子的关联函数沿 $\mathbb{R}^n$ 是拓扑的，沿 $\mathbb{C}^m$ 是全纯的。

% [Tutor Note: 推导补足：下降方程 (3.61) 的递归结构是：$d'O^{(k-1)}$ 给出「沿被平凡化方向的变化」，$(QO)^{(k)}$ 处理 $O$ 本身并非严格 $Q$ 闭的情形；
% 在 $Q$ 上同调中第二项消失，$O^{(k)}$ 变成 $d'$-闭的 $k$-形式，其积分只依赖同调类。因此 $Q$ 上同调中沿 $\mathbb{R}^n$ 的平移（拓扑性）与沿 $\bar z^{\bar a}$ 的平移（反全纯方向）都是平凡的，只剩 $z^a$ 的全纯依赖。]

\subsubsection{二维全纯扭变}

展现全纯下降的最简单理论，是二维 $\mathcal{N}=(0,2)$ 或 $\mathcal{N}=(2,2)$ 理论的全纯扭变 [20--22]；后者历史上称为半扭变理论。\footnote{尽管超出本讲义范围，近期有大量引人入胜的工作研究二维 $\mathcal{N}=(0,2)$ 理论的新颖对偶 [104, 105]，及其在四维流形几何中的应用 [106--108]。}
% [原文疑误: 原文 "are holomorphically twisted of 2d $\mathcal{N}=(0,2)$ or $\mathcal{N}=(2,2)$ theories" 缺 "versions of"，应为 "are holomorphically twisted versions of 2d $\mathcal{N}=(0,2)$ or $\mathcal{N}=(2,2)$ theories"；译文按后者理解。]
在此设定中有两个超荷 $Q_+,\bar Q_+$，其非平凡括号为 $\{\bar Q_+,Q_+\}=-2i\partial_{\bar z}$，于是例如 $Q=\bar Q_+$ 是一个全纯超荷，且 $Q_{\bar z}=\frac{i}{2}Q_+$ 平凡化 $\partial_{\bar z}$。
$Q$ 闭局域算子的关联函数只全纯地依赖于其位置，这一事实意味着碰撞赋予局域算子以顶点代数或手性代数的结构。

全纯扭变的二维 $\mathcal{N}=(0,2)$ 或 $\mathcal{N}=(2,2)$ 理论中的局域算子 $O$ 只有一个后代 $O^{(1)}=-d\bar z[Q_{\bar z},O]$，满足 $QO^{(1)}=\bar\partial O$。
注意，当 $QO=0$ 时，2-形式取值的局域算子 $dzO^{(1)}$ 的 $Q$ 变分是一个 1-形式取值的局域算子的外微分：$Q(dzO^{(1)})=d(dzO)$。
正是这个后代保证了（乘以 $dz$ 的）$Q$ 闭局域算子的围道积分只依赖于积分围道的同调类：若 $D$ 是 $\mathbb{C}$ 上某个紧连通区域，则
\begin{equation}\tag{3.62}
Q\int_D dzO^{(1)}=\int_D d(dzO)=\oint_{\partial D}dzO.
\end{equation}
% [Tutor Note: 推导补足：把 $Q$ 移入积分并用 $Q(dzO^{(1)})=d(dzO)$，再用 Stokes 定理；反向使用该式即得：围道在 $D$ 内连续形变不改变积分值。
% 这与拓扑情形中 $QO^{(k)}=dO^{(k-1)}$ 的论证完全平行，只是此处只允许全纯（复结构保持）的形变。]
当然，这种在形变积分围道下的不变性，从 $Q$ 上同调的全纯本性来看是意料之中的。

若有两个局域算子 $O_1$ 与 $O_2$ 分别位于点 $z$ 与 $w$，则可利用这一形变不变性来定义次级乘积\footnote{我们称之为「次级乘积」，但它更接近一维拓扑理论中的（分次）对易子。特别地，该括号次数为 $0$，并直接与初级乘积相联系；它只是提取 OPE 中的单极点。下文描述的 $\lambda$-括号可解释为 OPE 的傅里叶变换，参看 [109, 第 2 章]。}，例如把 $dzO_1$ 沿以 $w$ 为圆心的小圆积分。
更一般地，可通过与 $dze^{\lambda(z-w)}$ 作积分来定义它们的 $\lambda$-括号 $\{-\,_\lambda-\}$，它是位于 $w$ 的局域算子：
\begin{equation}\tag{3.63}
\{O_1{}_\lambda O_2\}(w)=\oint_{S^1}(dze^{\lambda(z-w)}O_1(z))O_2(w).
\end{equation}
% [Tutor Note: $\lambda$-括号把 OPE 的极点数据打包为 $\lambda$ 的多项式：$\{O_1{}_\lambda O_2\}=\sum_{n\ge0}\frac{\lambda^n}{n!}O_1{}_{(n)}O_2$，
% 其中 $O_1{}_{(n)}O_2$ 是 OPE 中 $(z-w)^{-n-1}$ 的系数；$\lambda^0$ 项即上文关于次级乘积的脚注中所说的单极点。]

% ==== 原文 PDF 第 57 页 ====
这是任何全纯场论中都存在的最简单的球代数（sphere algebras）实例之一，参看 [55, 110]。

在名副其实的全纯理论中，例如二维手征 CFT，$\lambda$-括号满足各种性质，所得结构有时称为 Lie 共形代数；与正规序乘积 $:O_1O_2:$ 合在一起，所得结构等价于通常的顶点代数概念 [111]。
为保持讨论简短，我们将主要关注该结构的 Lie 部分，即只涉及 $\lambda$-括号的那些性质。
% [背景补充: Lie 共形代数（Lie conformal algebra）是顶点代数「只保留 $\lambda$-括号」的版本：其公理为 sesquilinearity、分次反对称性与 $\lambda$-Jacobi 恒等式；补上正规序乘积后即得到完整顶点代数。]

在扭变理论中，Lie 共形代数结构的定义关系可能是自动成立的，也可能需要以同调方式施加。
例如，全纯导数 $\partial_w$ 的作用由下式给出：
\begin{equation}\tag{3.64}
\partial_w\{O_1{}_\lambda O_2\}(w)=\lambda\{O_1{}_\lambda O_2\}(w)+\{O_1{}_\lambda\partial O_2\}(w),
\end{equation}
其中 $\partial O_2(w)=\partial_wO_2(w)$。类似地有
\begin{equation}\tag{3.65}
\{\partial O_1{}_\lambda O_2\}(w)=-\lambda\{O_1{}_\lambda O_2\}(w).
\end{equation}
这些性质（合称为 $\lambda$-括号的半线性性，sesquilinearity）直接由 $\lambda$-括号的定义得到，并非以同调方式施加。

更有趣的是，$\lambda$-括号有分次反对称性的版本
\begin{equation}\tag{3.66}
\{O_1{}_\lambda O_2\}=-(-1)^{F(O_1)F(O_2)}\{O_2{}_{-\lambda-\partial}O_1\}.
\end{equation}
这通常如下验证：
\begin{equation}\tag{3.67}
\begin{aligned}
\{O_1{}_\lambda O_2\}(0)
&=\oint dz\,e^{\lambda z}O_1(z)O_2(0)
=\oint dz\,e^{\lambda z}\bigl(e^{z\partial+\bar z\bar\partial}O_1(0)\bigr)\bigl(e^{z\partial+\bar z\bar\partial}O_2(-z)\bigr)\\
&=(-1)^{F(O_1)F(O_2)}\oint dz\,e^{z(\lambda+\partial)}\bigl(O_2(-z)O_1(0)\bigr)\\
&=-(-1)^{F(O_1)F(O_2)}\{O_2{}_{-\lambda-\partial}O_1\}(0).
\end{aligned}
\end{equation}
% [Tutor Note: 推导补足：第二行把 $O_1(z)$、$O_2(0)$ 写成平移后的 $O_1(0)$、$O_2(-z)$，并把平移因子 $e^{\lambda z}$ 与 $e^{z\partial}$ 合并为 $e^{z(\lambda+\partial)}$；
% 第三行交换两个算子的次序，产生 Koszul 符号 $(-1)^{F(O_1)F(O_2)}$，并把 $\bar z\bar\partial$ 作用于 $O_2$ 的部分去掉（全纯性）；
% 最后一行把 $e^{-z\partial}O_2(-z)$ 按 Taylor 展开识别为 $\{O_2{}_{-\lambda-\partial}O_1\}$ 的定义，即得反对称性。]
注意，理论的全纯性被用来去掉平移 $e^{\bar z\bar\partial}$ 中的反全纯部分。
当这种全纯性以同调方式施加时，我们看到该反对称关系也必须以同调方式施加。
构造这样一个同伦的关键一步在于处理反全纯平移；特别地，我们需要构造一个同调次数为 $F(O)-1$ 的算子 $H_{\bar z}(O(0))$，使得：
\begin{equation}\tag{3.68}
O(z)\bigl(\equiv e^{z\partial+\bar z\bar\partial}O(0)\bigr)=e^{z\partial}O(0)+QH_{\bar z}(O(0)).
\end{equation}
利用下降超荷 $Q_{\bar z}$ 可以直截了当地构造这样的同伦，例如
\begin{equation}\tag{3.69}
\begin{aligned}
H_{\bar z}(O(0))
&=e^{z\partial}\int_0^{\bar z}d\bar z'\,e^{\bar z'\bar\partial}O(0)
=\sum_{n\ge0,\ \bar n\ge1}\frac{z^n\bar z^{\bar n}}{n!\bar n!}\partial^n\bar\partial^{\bar n-1}Q_{\bar z}O(0),\\
&\rightsquigarrow QH_{\bar z}(O(0))
=\sum_{n\ge0,\ \bar n\ge1}\frac{z^n\bar z^{\bar n}}{n!\bar n!}\partial^n\bar\partial^{\bar n}O(0)
=O(z)-e^{z\partial}O(0).
\end{aligned}
\end{equation}
% [Tutor Note: 推导补足：$H_{\bar z}$ 是对 $\bar z$ 积分（从 $0$ 到 $\bar z$）得到的「平移同伦」；逐项求导给出 $\bar\partial^{\bar n}$，再用 $\{Q,Q_{\bar z}\}=iP_{\bar z}\sim i\bar\partial$（见 (3.60)）
% 把 $Q_{\bar z}O$ 换成 $\bar\partial O$，即得最后一行：$QH_{\bar z}$ 恰是 $e^{\bar z\bar\partial}$ 展开中除去 $e^{z\partial}$ 之后的部分。]

% ==== 原文 PDF 第 58 页 ====
由于我们处理的是单个复维度，平移没有更高阶的同伦——那将需要 $p>1$ 的时空 $(0,p)$ 形式。
我们可以把这一平移同伦在复平面上的 2 点构型空间 $C_{\mathbb{C}}(2)$ 上积分，从而构造（分次）反对称性所需的同伦。\footnote{我们用 $\mathbb{C}$ 而非 $\mathbb{R}^2$ 来强调该构型空间的复本性，即它具有从 $\mathbb{C}$ 继承的自然复结构。}
例如，记
\begin{equation}\tag{3.70}
H(O_1,O_2)(0)=\oint dz\,e^{\lambda z}H_{\bar z}(O_1(0)O_2(-z));
\end{equation}
若 $O_1$ 与 $O_2$ 是 $Q$ 闭的，则
\begin{equation}\tag{3.71}
QH(O_1,O_2)(0)=\{O_1{}_\lambda O_2\}(0)+(-1)^{F(O_1)F(O_2)}\{O_2{}_{-\lambda-\partial}O_1\}(0).
\end{equation}
% [Tutor Note: (3.71) 表明括号的分次反对称性的失效量恰好是 $Q$ 恰当的，同伦由 (3.70) 的 $H(O_1,O_2)$ 给出；
% 其同调次数为 $F(O_1)+F(O_2)-1$，故 $QH$ 与 $\{O_1{}_\lambda O_2\}$ 的次数相同。]

更有趣的是 $\lambda$-Jacobi 恒等式：对任意局域算子 $O_1,O_2,O_3$ 与 $\lambda,\mu\in\mathbb{C}$，该括号满足
\begin{equation}\tag{3.72}
\{O_1{}_\lambda\{O_2{}_\mu O_3\}\}=\{\{O_1{}_\lambda O_2\}_{\lambda+\mu}O_3\}+(-1)^{F(O_1)F(O_2)}\{O_2{}_\mu\{O_1{}_\lambda O_3\}\}.
\end{equation}
正如括号的分次反对称性一样，这一关系一般以同调方式施加。
我们可以把所需同伦在拓扑设定中三元 $L_\infty$ 运算 $\ell_3$ 所用的同一区域 $\mathcal{L}_3$（见 Figure 14）上积分来构造它：
\begin{equation}\tag{3.73}
\{O_1{}_\lambda O_2{}_\mu O_3\}:=\int_{\mathcal{L}_3}(du\,e^{\lambda(u-w)}O^{(1)}_1)(dz\,e^{\mu(z-w)}O_2)O_3.
\end{equation}
由于这一运算涉及单个后代，可知它是次数 $-1$ 的运算。
直接应用全纯下降方程与 Stokes 定理可推出：对 $Q$ 闭的 $O_i$，$\{O_1{}_\lambda O_2{}_\mu O_3\}$ 的 $Q$ 变分恰好是 $\lambda$-Jacobi 子：
\begin{equation}\tag{3.74}
Q\{O_1{}_\lambda O_2{}_\mu O_3\}=\{O_1{}_\lambda\{O_2{}_\mu O_3\}\}-\{\{O_1{}_\lambda O_2\}_{\lambda+\mu}O_3\}-(-1)^{F(O_1)F(O_2)}\{O_2{}_\mu\{O_1{}_\lambda O_3\}\},
\end{equation}
其中右端第一项来自 $D$ 的外边界，第二项来自 $D$ 中围绕 $w$ 的边界，第三项来自 $D$ 中围绕 $z$ 的边界。
% [Tutor Note: 推导补足：在 $\mathcal{L}_3$ 上积分并把 $Q$ 移入被积函数，$Q$ 作用在 $O^{(1)}_1$ 上产生 $\bar\partial O_1$，于是 Stokes 定理给出三块边界：
% 外边界（对应 (3.74) 第一项）、围绕 $w$ 的小圆（第二项）、围绕 $z$ 的小圆（第三项）。这与拓扑情形中 $\ell_3$ 给出 Jacobiator 的论证完全平行，
% 区别只在于此处积分的区域是复构型空间中的实三维链 $\mathcal{L}_3$。]
毫不意外，存在整座更高 $n$ 元括号的塔，它们充当更高 Jacobi 恒等式之间的同伦，与 $L_\infty$ 代数完全类似；我们称所得结构为 $L_\infty$ 共形代数。\footnote{近期论文 [112] 引入了 $n$-Lie 共形代数这一相关概念。该结构对应于我们所描述的 $L_\infty$ 共形代数，其中只有 $n$ 元括号 $\{O_1{}_{\lambda_1}\dots{}_{\lambda_{n-1}}O_n\}$ 非零。}
一旦与正规序乘积以及实现定义关系 $Q$ 恰当性的各种同伦相结合，这一 $L_\infty$ 共形代数结构应当（类比于无同伦的设定）延拓为同伦顶点代数。\footnote{更细致地研究这些同伦，并把所得结构与 Francis--Gaitsgory [98] 发展的同伦手征代数概念相比较，将是很有意义的。}

% ==== 原文 PDF 第 59 页 ====
\subsubsection{三维全纯–拓扑扭变}

混合全纯–拓扑理论中的下降是拓扑理论下降与全纯理论下降的耐人寻味的混合。
许多简单例子来自具有（至少）$\mathcal{N}=2$ 超对称的三维理论。考虑三维 $\mathcal{N}=2$ 超代数：
\begin{equation}\tag{3.75}
\{Q_+,\bar Q_+\}=-2i\partial_{\bar z} \qquad \{Q_-,\bar Q_-\}=2i\partial_z \qquad \{Q_-,\bar Q_+\}=\{Q_+,\bar Q_-\}=i\partial_t.
\end{equation}
由此，例如 $Q:=\bar Q_+$ 是一个全纯–拓扑超荷，且导数 $\partial_t,\partial_{\bar z}$ 通过 $Q_t=-iQ_-$ 与 $Q_{\bar z}=\frac{i}{2}Q_+$ 而为 $Q$ 恰当的。
% [Tutor Note: 物理图像：三维 $\mathcal{N}=2$ 的 $R$ 对称群为 $U(1)$；取 $Q=\bar Q_+$ 后，$Q$ 只平凡化 $t$ 与 $\bar z$ 方向的平移（$Q_t,Q_{\bar z}$ 存在），
% 故扭变理论沿 $t$ 拓扑、沿 $z$ 全纯——这正是「全纯–拓扑」名称的来源。]
与二维一样，碰撞赋予 $Q$ 闭局域算子以顶点代数的结构。
重要的是，两个 $Q$ 闭局域算子的 OPE 必然非奇异（模去 $Q$ 恰当项）：这些 OPE 只可能在插入点重合的极限下奇异，但同时（模去 $Q$ 恰当项）与 $t$ 方向上的间距无关，因此在 $Q$ 上同调中非奇异。
所得顶点代数称为交换的。\footnote{这个名字的理由在于：非奇异的 OPE 蕴含模式（modes）构成的 Lie 代数是阿贝尔的，即（分次）交换的。}
这是混合全纯–拓扑理论的一般特征。

扭变的三维 $\mathcal{N}=2$ 理论也容许 $\lambda$-括号，此时它由将 $dze^{\lambda(z-w)}O^{(1)}_1(t,z,\bar z)$ 在围绕 $O_2(s,w,\bar w)$ 的 $S^2$ 上积分给出：
\begin{equation}\tag{3.76}
\{O_1{}_\lambda O_2\}(s,w,\bar w)=\int_{S^2}(dz\,e^{\lambda(z-w)}O^{(1)}_1(t,z,\bar z))O_2(s,w,\bar w).
\end{equation}
由于这一运算涉及 1 个后代，$\lambda$-括号现在的次数为 $-1$。
与（交换的）OPE 合在一起，所得结构可称为 $(-1)$-平移的 Poisson 顶点代数。
于是我们得出结论：全纯扭变的三维理论中局域算子的 $Q$ 上同调自然地具有这一结构 [56, 74]。
% [原文疑误: 原文 "the $Q$-cohomology ... is naturally have this structure" 语法不通，应为 "naturally has this structure"（或 "is naturally given this structure"）；依据为动词单复数搭配，译文按 "has" 处理。]
% [Tutor Note: 物理图像：三维全纯–拓扑扭变的 $\lambda$-括号次数为 $-1$，OPE 交换；这正是 Poisson 顶点代数的「平移」版本。
% 对一般的 $k$-平移 Poisson 顶点代数，$\lambda$-括号的次数为 $k$；这里 $k=-1$ 来自只用了单个后代（一次下降）。]

与分别出现在拓扑一维与二维理论中的 $A_\infty$ 代数与 $L_\infty$ 代数一样，在链层面工作应当能给出关于一般混合全纯–拓扑理论中局域算子的碰撞与下降的更精细信息。
例如，$(-1)$-平移的 $\lambda$-括号应当满足 $(-1)$-平移版本的 $\lambda$-Jacobi 恒等式，后者如二维情形那样由积分后代的 $Q$ 变分以同调方式施加。
类似地，拓扑 $t$ 方向上的下降应当容许该 $\lambda$-括号的 $A_\infty$ 丰富化。
关于这种被他们称为「$A_\infty$-手征代数」的结构会包含什么内容，更多细节见 [56, 第 2.4 节]。

同样值得一提的是，第 2.4.1 节末尾所述的 Hochschild 上同调构造有一个适用于这一混合三维设定的类比。
特别地，全纯–拓扑扭变中体局域算子的代数应当可以从一个（足够大的）边界条件恢复。
任何这样的边界条件都容许一个局域算子的（同伦）顶点代数，而体局域算子的代数可以从其（导出的）中心提取，参看 [56, 第 2.4 节]。
Zeng 近期的工作 [113] 利用这一构造描述了扭变的三维 $\mathcal{N}=2$ 阿贝尔规范理论中局域算子的交换顶点代数的模式代数（mode algebra）的若干侧面；该代数因单极子算子而获得有趣的非微扰修正。

% ==== 原文 PDF 第 60 页 ====
\subsubsection{顶点代数与 Poisson 顶点代数的同伦作用}

自然可以预期，出现在扭变超对称场论与弦论中的顶点代数与（平移的）Poisson 顶点代数的实例，具有涉及上述同伦代数的链层面提升。
详细理解这些顶点代数的同伦丰富化如何在对偶下传递，将是令人欣慰的。
有许多美妙的顶点代数例子（例如可从扭变理论得到的那些）在对偶映射下传递。
弦论中的一个经典例子出现在四环面 $T^4$ 上的杂化弦（heterotic string）与 K3 面上的 type IIA 弦的对偶中。
杂化弦中的 BPS 态与世界面上的顶点算子相关，使得超对称的右行模式处于其基态；在半扭变中，人们保留一个全纯顶点代数结构，它本质上来自左行模式对 BPS 态的作用。
杂化弦的 BPS 态在 type IIA 中映射为 K3 上各种包裹 D 膜的构型。
更确切地说，BPS 态由 K3 上半稳定凝聚层模空间的同调中的微分形式表示；细节见 [114]。
在杂化弦一侧显然的代数作用，以对应的方式几何地作用在这些瞬子模空间上，其方式最早由 Nakajima [115, 116] 与 Grojnowski [117] 阐明。

近期有许多美妙的顶点代数例子，它们作用在数学家和物理学家感兴趣的各种模空间上。
例如，Gaiotto--Rapcak 的工作识别出一大类「角顶点代数（corner VOAs）」$Y_{(L,M,N)}[\Psi]$ [118]，它们出现在由 IIB 弦论中基本 5-膜汇合所诱导的扭曲 $\mathcal{N}=4$ 超 Yang--Mills 理论的 $\frac{1}{2}$-BPS 边界条件的界面上。
弦对偶把这一设定与 type IIA 弦论在 Calabi--Yau 三维流形上包裹 D 膜的构型联系起来，而这些 VOA 也被期望描述 Omega 形变 M 理论中某一组 M5 膜上的局域算子代数 [119]。
角 VOA [120] 作用于 Nekrasov 发现的尖刺瞬子（spiked instantons）模空间的等变上同调 [121--126]；反过来，这可以看作 $W$ 代数 $W^\Psi_N=Y_{(0,0,N)}[\Psi]$ 对 $U(N)$ 瞬子模空间（的等变上同调）的作用的推广，后者由 AGT 猜想 [129] 预言 [127, 128]。
由包裹四维流形（即 Calabi--Yau 三维流形中的除子）的 D0 与 D2 膜构型得到的 BPS 代数，也已从 4d--2d 对应的角度得到研究 [106, 107, 130]。
来自带边界三维 $\mathcal{N}=2$ 理论的全纯–拓扑扭变的顶点代数 [45, 56, 106]，与上述弦论及 M 理论构造相关 [105, 106, 108, 131]，并且本身也已成为许多美妙洞见的近期来源，例如 [132--134]。
% [背景补充: 「角 VOA」得名于它们出现在四维理论中两个（半）BPS 边界条件相交的「角」处：两个边界条件的交线（codimension-2 界面）承载一个顶点代数，是普通边界 VOA 的推广。]

这一推理路线提出了（目前尚无答案的）问题：这些（角与边界）VOA 的链层面提升是什么？所得的（同伦）VOA 是否作用在相应的模空间同调的（适当的链层面实现）上？\footnote{这些进展还触及物理与数学中众多深刻结构，例如 BPS 态代数与上同调 Hall 代数、计数几何中的 Donaldson--Thomas 理论、解析数论中的（mock 与量子）模性、低维流形的拓扑等等，远远超出本讲义的范围。鼓励读者研读所引文献及其中的参考文献。}
